Fix \(0<\Delta\leq\Delta_0\).  By \(\mathrm{lem:boundary\text{-}anchor\text{-}uniqueness}\), \(\tau_P(b)\) is a well-defined functional for every \(P\in\mathcal P_L\).

First fix \(Q\in\mathcal Q_\Delta\).  By \(\mathrm{def:released\text{-}identified\text{-}set}\), there is a law \(P\in\mathcal P_L\) satisfying \(\mathcal L_P(Z_\Delta)=Q\).  For \(t\in\{0,1\}\), define the released event
\[
 A_{t,\Delta}=\{T=t,Q_\Delta(R)=0\}.
\]
Since \(R\in[0,1]\), \(\Delta>0\), and \(Q_\Delta(r)=\Delta\lfloor r/\Delta\rfloor\),
\[
 A_{t,\Delta}=\{T=t,0\leq R<\Delta\}=\{T=t,R<\Delta\}.
\]
The two-sided local-mass condition in \(\mathrm{ass:uniform\text{-}disc\text{-}design}\), together with \(0<\Delta\leq\Delta_0\leq1/4<1/2\), therefore gives
\[
 Q(A_{t,\Delta})
 =\mathbb P_P(T=t,R<\Delta)
 \geq c_0\Delta^2>0.
\]
This probability depends only on the released law.  Hence every attainable \(Q\) assigns positive probability to both first-bin side events, and the released-law conditional mean
\[
 \nu_t(Q)=\mathbb E_Q[Y\mid A_{t,\Delta}]
\]
is well defined.

Let \(P'\in\mathcal P_L\) be any other member of the same fiber, so that \(\mathcal L_{P'}(Z_\Delta)=Q\).  Write \(\mu'_t\) and \(\theta'_t\) for its side regression and boundary anchor.  Equality of the released laws makes \(Q(A_{t,\Delta})\) and \(\nu_t(Q)\) common to \(P\) and \(P'\).  Apply the general-event identity in \(\mathrm{lem:conditional\text{-}local\text{-}binomial\text{-}law}\), with its auxiliary sample-size argument fixed at the admissible value \(n=256\), to \(B=A_{t,\Delta}\) under each law.  The event belongs to \(\sigma(T,Q_\Delta(R))\), is contained in \(\{T=t\}\), and has the strictly positive probability just proved, so all premises of that identity hold.  It follows that
\[
 \nu_t(Q)
 =\mathbb E_P[\mu_t(X)\mid A_{t,\Delta}]
 =\mathbb E_{P'}[\mu'_t(X)\mid A_{t,\Delta}].
\]
On \(A_{t,\Delta}\), \(R<\Delta\leq1/2\).  Thus the point-anchor envelope in \(\mathrm{ass:radial\text{-}lipschitz}\), the conditional triangle inequality, and positivity of the conditioning event yield
\[
 \begin{aligned}
 |\theta_t-\nu_t(Q)|
 &=\left|\mathbb E_P[\theta_t-\mu_t(X)\mid A_{t,\Delta}]\right|\\
 &\leq\mathbb E_P[|\theta_t-\mu_t(X)|\mid A_{t,\Delta}]\\
 &\leq L\mathbb E_P[R\mid A_{t,\Delta}]\\
 &\leq L\Delta,
 \end{aligned}
\]
and, by the same calculation under \(P'\),
\[
 |\theta'_t-\nu_t(Q)|\leq L\Delta.
\]
Consequently, for each \(t\in\{0,1\}\),
\[
 |\theta_t-\theta'_t|
 \leq|\theta_t-\nu_t(Q)|+|\theta'_t-\nu_t(Q)|
 \leq2L\Delta.
\]
Using the definition \(\tau_P(b)=\theta_1-\theta_0\), we obtain
\[
 \begin{aligned}
 |\tau_P(b)-\tau_{P'}(b)|
 &=| (\theta_1-\theta_0)-(\theta'_1-\theta'_0)|\\
 &\leq|\theta_1-\theta'_1|+|\theta_0-\theta'_0|\\
 &\leq4L\Delta.
 \end{aligned}
\]
Every pair of elements of \(I_\Delta(Q)\) arises from such a pair \(P,P'\) by \(\mathrm{def:released\text{-}identified\text{-}set}\).  Taking the supremum over pairs in this fiber and then over \(Q\in\mathcal Q_\Delta\) proves
\[
 \omega_\Delta\leq4L\Delta.
\]

For the reverse inequality, \(\mathrm{lem:first\text{-}bin\text{-}alias}\) supplies laws \(P^{\mathrm{alias}}_{+,\Delta},P^{\mathrm{alias}}_{-,\Delta}\in\mathcal P_L\) having a common one-unit released law, denoted by \(Q^{\mathrm{alias}}_\Delta\), and satisfying
\[
 \tau_{P^{\mathrm{alias}}_{+,\Delta}}(b)
 -\tau_{P^{\mathrm{alias}}_{-,\Delta}}(b)
 =2c_L\Delta.
\]
Their class membership implies \(Q^{\mathrm{alias}}_\Delta\in\mathcal Q_\Delta\), and both target values belong to \(I_\Delta(Q^{\mathrm{alias}}_\Delta)\).  Therefore
\[
 \omega_\Delta
 \geq\operatorname{diam}I_\Delta(Q^{\mathrm{alias}}_\Delta)
 \geq2c_L\Delta.
\]
Because \(c_L>0\) and \(\Delta>0\), this attainable released law is compatible with two distinct values of the full-data causal target.  Hence \(\tau_P(b)\) is generally not point identified from a fixed-\(\Delta\) release.  The bounds
\[
 2c_L\Delta\leq\omega_\Delta\leq4L\Delta
\]
also show that its worst release-fiber diameter is of order \(\Delta\), with constants independent of \(\Delta\).